\section{Background and Related Work}

FinOps has matured into an operational framework for technology value management that emphasizes collaboration between engineering, finance, and business stakeholders, with iterative phases commonly described as Inform, Optimize, and Operate \cite{finops_framework,finops_phases}. That framing is useful for this paper because cost management is not merely retrospective reporting. It requires timely data, attribution structures, optimization logic, and a feedback path from proposed action to realized outcome.

\subsection{Multi-Cloud Cost and Attribution Context}
Prior work on chargeback for cloud services argues that accountability, transparency, and measurability are central to successful cloud cost allocation \cite{baars2014chargeback}. Those concerns become sharper in multi-cloud settings, where billing and SLA semantics differ across providers and make manual comparison difficult \cite{gupta2016streamlining,imran2020multicloud}. In practice, multi-cloud consumers must reconcile provider-specific exports, normalize ownership and pricing fields, and decide how to treat shared or unattributed spend before any optimization recommendation can be trusted.

\subsection{Provider-Native Billing Semantics}
The major providers already expose rich billing exports, but they do so with different semantics. AWS Cost and Usage Reports distinguish line-item types such as \texttt{Usage}, \texttt{DiscountedUsage}, \texttt{SavingsPlanCoveredUsage}, \texttt{RIFee}, and \texttt{SavingsPlanRecurringFee}, and also provide reservation-effective-cost columns that are populated only for certain row types \cite{aws_line_items,aws_reservation_details}. Azure Cost Management distinguishes actual and amortized cost views and surfaces explicit unused commitment charge types such as \texttt{UnusedReservation} and \texttt{UnusedSavingsPlan} only in amortized analysis \cite{azure_amortized_cost,azure_unused_reservation}. Google Cloud Billing exports attach discounts and credits directly to usage records and document how committed-use discounts appear through commitment fees, effective prices, and credit offsets \cite{gcp_cud_overview,gcp_standard_export}. These native exports are necessary inputs, but they are not themselves a cross-cloud analytical contract.

\subsection{Standards and Differentiation}
The FinOps Open Cost and Usage Specification (FOCUS) is an important standardization effort for cross-provider billing data \cite{focus_v12}. It reduces naming friction and encourages a common schema vocabulary, but it does not by itself define cloud-specific NEC formulas, shared-cost allocation behavior, recommendation logic, or lifecycle validation outputs. Existing literature also contains broad surveys of cloud cost optimization strategies \cite{liu2023cost}, yet much of that work focuses on provider-local optimization rather than on an implemented end-to-end decision pipeline for post-billing FinOps.

This paper therefore positions its contribution between provider-native documentation and abstract optimization literature. Relative to native tooling, it contributes an open and reproducible cross-cloud decision substrate. Relative to billing standards, it contributes explicit NEC, allocation, and validation semantics. Relative to prior research on chargeback and multi-cloud management, it contributes an implementation-backed pipeline that links normalization, attribution, signals, recommendations, forecasting, and lifecycle-aware evaluation in one auditable system.
